% Source: https://artofproblemsolving.com/community/c3h1734443_geometry_questions_from_a_local_math_competition
% Stable identity: aops-topic-1734443-q08
A regular dodecahedron has $12$ faces, each a regular pentagon with side $10$. Split a face into $5$ isosceles triangles from its centre $O$; each has apex angle $72^\circ$, base $10$ and height equal to the apothem $a$:

\begin{asy}
usepackage("amsmath");
size(4.5cm);
defaultpen(fontsize(10pt));
real s=10; real ap=(s/2)/Tan(36); pair O=(0,0);
pair[] P; for(int i=0;i<5;++i) P.push((s/(2*Sin(36)))*dir(270-36+72*i));
draw(P[0]--P[1]--P[2]--P[3]--P[4]--cycle,linewidth(0.9));
pair M=(P[0]+P[1])/2;
draw(O--M,blue+dashed);draw(O--P[0],gray(0.5));draw(O--P[1],gray(0.5));
dot(O);label("$O$",O,N);
label("$10$",M,S);label("$a$",(O+M)/2,E,blue);
label("$36^\circ$",O+3.8*dir(252),fontsize(7pt));draw(arc(O,2.0,234,270));
label("\mbox{}",(min(currentpicture,user=true)+max(currentpicture,user=true))/2,black); // keeps the text after the figure black when LaTeX typesets the labels inline
\end{asy}

Since half of the apex angle is $36^\circ$, $a=\dfrac{5}{\tan36^\circ}$, and the face area is $5\cdot\dfrac{1}{2}\cdot10\cdot a=\dfrac{125}{\tan36^\circ}$. Using $\tan36^\circ=\sqrt{5-2\sqrt5}$, this equals $25\sqrt{25+10\sqrt5}$, since $\dfrac{125}{\sqrt{5-2\sqrt5}}=\dfrac{125\sqrt{5+2\sqrt5}}{\sqrt{25-20}}=25\sqrt5\sqrt{5+2\sqrt5}$. So the surface area is
\[
12\cdot25\sqrt{25+10\sqrt5}=\boxed{300\sqrt{25+10\sqrt5}}\approx2064.7 .
\]
